%%==================================================
%% thanks.tex for BIT Master Thesis
%% modified by yang yating
%% version: 0.1
%% last update: Dec 25th, 2016
%%==================================================

\begin{thanks}

时光荏苒，硕士研究生的学习生涯即将画上句号。回首这段充满挑战与收获的岁月，心中涌动着无限的感激之情。在此，谨向所有关心、支持和帮助过我的师长、同学、亲友致以最诚挚的谢意。

古之学者必有师。师者，所以传道授业解惑也。感谢我的导师周猛老师，对我科研思维的训练，以及让我全方面认识了自己课题。周老师在科研方面有自己的坚持与热爱，并且他按照自己的方法引导每一个学生对自己课题的思考，同时也给予了我们非常轻松的学习环境。老师在选题，写作，修改等环节的认真态度，以及在科研方面的严谨也深深影响到了我。从研一“看山是山”的主体，到研二“看山不是山”的客体，最后又回归研三“看山还是山”的本我主体，这一路正因为有周老师的指导，才让我快速领悟到了科研的魅力，以及启发自己对所做课题现实意义和个人未来在社会存在方式的思考。三年地靠近课题，靠近自己，靠近“花开随喜，花落不悲”，让我找到了“独而近己”背后的真正意义。周老师曾说“什么样的年龄就要干什么样的事”，这句话我记忆犹新。万物在不同的阶段，有着不同的位置，思想以及信息差和追求，而我们就要有自己在所处阶段该有的样子，做自己可控范围的事儿。

最后，感谢父母对我的养育之恩，以及给予我的温室，是在他们的鼓励和支持下一直前行着。

\hfill 2026年5月

\hfill 李世龙
\end{thanks}